\chapter{Reprendre après une interruption}

Sur dix jours d'entraînement, une coupure n'est pas une éventualité : c'est une
certitude. Panne de courant, mise à jour qui redémarre la machine, besoin
soudain de la carte graphique pour autre chose. Ce chapitre explique comment
tout cela devient sans conséquence.

\section{La bonne nouvelle, tout de suite}

\begin{nkcode}[]{Reprendre : la même commande, exactement}
NKQwen2Train --gguf=<modele> --corpus=<corpus.txt> --out=mon-modele
\end{nkcode}

C'est tout. Aucune option de reprise à ajouter. Le programme cherche les
checkpoints, prend le plus avancé qui soit intact, et repart.

\begin{nkcode}[]{Ce qu'il affiche alors}
REPRISE : pas 6, exemple 6/100017, époque 0
\end{nkcode}

\begin{nkretenir}
Relancer la même commande suffit. Si vous voyez « REPRISE », le travail
précédent est retrouvé. Si vous voyez « aucun checkpoint exploitable », il
repart de zéro — et il faut comprendre pourquoi avant de le laisser tourner.
\end{nkretenir}

\section{Pourquoi recharger les poids ne suffit pas}

Voici l'idée la plus subtile du cours. Prenez le temps.

Naïvement, on croirait qu'il suffit de sauvegarder les paramètres appris. C'est
faux, et une reprise faite ainsi \textbf{diverge dès le premier pas}.

La raison tient à la méthode qui décide de l'ampleur des corrections, appelée
\textbf{Adam}. Adam ne se contente pas du gradient courant : il garde une
mémoire de ceux d'avant, sous la forme de deux moyennes appelées
\emph{moments}. C'est ce qui lui permet d'avancer avec régularité plutôt que de
zigzaguer.

Or ces moyennes démarrent à zéro, donc sont trop faibles au début. Adam corrige
ce défaut par un facteur qui dépend du \textbf{numéro du pas} : très fort au
pas~1, négligeable au pas~1000.

Maintenant, imaginez une reprise qui recharge les poids mais oublie le
compteur :

\begin{itemize}
  \item Les moments sont grands (mille pas d'histoire) ;
  \item mais le compteur est retombé à zéro, donc Adam applique la correction
        du \emph{tout premier} pas ;
  \item il en résulte un pas gigantesque, et l'entraînement part en vrille.
\end{itemize}

Dix jours de calcul détruits en une correction.

\begin{nkretenir}
Un checkpoint doit contenir \textbf{quatre} choses, pas une : les paramètres,
les deux moments d'Adam, le numéro du pas, et la position dans le corpus. C'est
la différence entre reprendre et recommencer en pire.
\end{nkretenir}

\section{Ce que contient réellement un checkpoint}

\begin{center}
\begin{tabular}{p{4.6cm}p{9.0cm}}
\toprule
Les paramètres LoRA & Ce que le modèle a appris. \\
Les deux moments d'Adam & La mémoire de l'optimiseur — deux fois la taille des
                          paramètres. \\
Le numéro du pas & Pour que la correction de biais reste juste. \\
La position dans le corpus & Pour ne pas réapprendre les mêmes exemples. \\
Le numéro d'époque & Combien de fois le corpus a été parcouru. \\
L'état du tirage aléatoire & Pour que l'ordre reste reproductible. \\
Le taux d'apprentissage et les réglages d'Adam & Pour reprendre à l'identique. \\
\bottomrule
\end{tabular}
\end{center}

D'où la taille : \textbf{231~Mo} contre 77~Mo pour les seuls adaptateurs. Les
moments pèsent deux fois les paramètres.

\section{Écrire sans jamais rien détruire}

Un checkpoint de 231~Mo met environ une seconde à s'écrire. Que se passe-t-il si
le courant saute pendant cette seconde ?

Si l'on écrivait directement sur le fichier existant : on perdrait à la fois
l'ancien (déjà écrasé) et le nouveau (incomplet). \textbf{La sauvegarde aurait
détruit ce qu'elle devait protéger.}

Deux protections travaillent ensemble.

\subsection{L'écriture passe par un fichier temporaire}

Le programme écrit dans \texttt{mon-modele\_ckpt0.nkla.tmp}, puis — une fois le
fichier complet — le \textbf{renomme}. Le renommage est quasi instantané,
là où l'écriture dure une seconde.

\subsection{Deux checkpoints en alternance}

Les sauvegardes vont tour à tour dans \texttt{\_ckpt0} et \texttt{\_ckpt1}. À
tout instant, l'un des deux est une génération complète et saine.

\begin{nkretenir}
Fichier temporaire puis renommage, et deux fichiers en alternance. Ces deux
mesures ensemble font qu'aucune coupure, à aucun moment, ne peut vous laisser
sans checkpoint valide.
\end{nkretenir}

\section{Quand un fichier s'abîme}

Un disque peut abîmer un fichier sans prévenir. Chaque checkpoint porte donc une
\textbf{empreinte} calculée sur l'ensemble de son contenu. Au chargement, elle
est recalculée et comparée.

Voici une expérience réelle. Un checkpoint a été délibérément abîmé — 64 octets
écrasés en plein milieu, exactement ce que ferait un secteur défectueux — puis
l'entraînement relancé :

\begin{nkcode}[]{Sortie réelle après corruption volontaire}
checkpoint mon-modele_ckpt0.nkla ignoré :
  NkLoraGpuLoadNKLA : empreinte FNV incorrecte — fichier corrompu
REPRISE : pas 10, exemple 10/100017, époque 0
\end{nkcode}

Corruption détectée, fichier écarté, reprise depuis l'autre — \textbf{sans
perdre un seul pas}.

\begin{nkpiege}
Un fichier corrompu est \textbf{refusé}, jamais chargé « au mieux ». Charger des
paramètres à moitié faux donnerait un modèle silencieusement dégradé : il
répondrait, mais mal, et vous chercheriez pendant des jours une erreur qui
serait en réalité une corruption de disque.
\end{nkpiege}

\section{Arrêter proprement}

Pour interrompre volontairement : \texttt{Ctrl+C}.

Le programme ne s'arrête pas net. Il termine le pas en cours, écrit un
checkpoint, puis quitte :

\begin{nkcode}[]{Arrêt volontaire}
interruption demandée — écriture de l'état avant de quitter
checkpoint final : mon-modele_ckpt1.nkla (231.0 Mo)
adaptateurs livrables : mon-modele.nkla (77.0 Mo, 20185088 paramètres)
\end{nkcode}

Vous récupérez donc la carte graphique immédiatement, sans rien perdre.

\section{Régler la fréquence des sauvegardes}

\begin{center}
\begin{tabular}{lp{9.6cm}}
\toprule
\texttt{-{}-save-every=50} & Sauvegarde fréquente. Vous perdez au plus 50 pas
   ($\approx$ 8 min), mais vous écrivez 231~Mo toutes les 8 minutes. \\
\texttt{-{}-save-every=200} & \textbf{Le défaut.} Environ toutes les 30 minutes.
   Bon compromis. \\
\texttt{-{}-save-every=1000} & Peu d'écritures, mais une coupure peut coûter
   2~h~30 de calcul. \\
\bottomrule
\end{tabular}
\end{center}

\begin{nkretenir}
Le bon réglage répond à une seule question : \emph{combien de temps de calcul
acceptez-vous de perdre ?} Une sauvegarde coûte une seconde ; une reprise
manquée coûte des heures.
\end{nkretenir}

\section{Repartir vraiment de zéro}

Si vous voulez ignorer tout ce qui existe et recommencer :

\begin{nkcode}[]{Ignorer les checkpoints}
NKQwen2Train --gguf=<modele> --corpus=<corpus.txt> --out=mon-modele --fresh
\end{nkcode}

\begin{nkpiege}
\texttt{-{}-fresh} n'efface pas les checkpoints : il les ignore. Ils seront
écrasés à la première sauvegarde du nouvel entraînement. Si vous teniez à
l'ancien, \textbf{copiez-le ailleurs avant}, ou changez de \texttt{-{}-out}.
\end{nkpiege}

\begin{nkexo}
Lancez un entraînement avec \texttt{-{}-steps=10 -{}-save-every=5}. Laissez-le
finir. Relancez la même commande : vous devez lire « REPRISE : pas 10 ».
Puis relancez avec \texttt{-{}-fresh} : vous devez lire « les checkpoints
existants sont ignorés ». Ces deux messages sont ceux que vous lirez pendant
des jours — apprenez à les distinguer d'un coup d'œil.
\end{nkexo}
